\documentclass[10pt]{article}

\usepackage{geometry}
\geometry{
	a4paper,
	total={6.85in, 9.92in},
	left=0.71in,
	top=0.63in,
}
\usepackage[utf8]{inputenc}
\usepackage{hyperref}
\usepackage{listings}
\usepackage{xcolor}

\pagenumbering{gobble}
\setlength{\parskip}{0.35em}
\setlength{\parindent}{0pt}

\definecolor{cmt}{RGB}{110,110,110}
\lstset{
	basicstyle=\ttfamily\small,
	commentstyle=\color{cmt}\itshape,
	keywordstyle=\bfseries,
	columns=fullflexible,
	frame=leftline,
	framesep=6pt,
	xleftmargin=8pt,
	showstringspaces=false,
	breaklines=true,
}

\title{Feeding Frenzy 2: unpacking a self-checksumming fish tank}
\author{}

\makeatletter
\newcommand{\fsize}{\f@size pt }
\renewcommand{\maketitle}{
\begin{flushleft}
{\noindent\Huge\bf\@title}\break
\end{flushleft}
}
\makeatother

\begin{document}
\maketitle

\bigskip

PopCap's \textit{Feeding Frenzy 2} (2006) keeps every sprite, sound, and
cutscene in one $\sim$14\,MB file: \texttt{FF2.saf}. No loose assets, just an
opaque blob. Here is how it comes apart, and the single trick that makes
putting it back together interesting.

The header is trivial:

\begin{lstlisting}[language=C]
struct saf_header {  // little-endian
    char magic[4];   // "FFAS"
    u32  version;    // 1
    u32  toc_offset; // jump here for index
};
\end{lstlisting}

Everything from \texttt{0x0C} to \texttt{toc\_offset} is the raw concatenation
of every packed file. The table of contents sits at the end:

\begin{lstlisting}[language=C]
struct toc {
    u32 version;
    u8  global_checksum[16];
    u32 entry_count;
    struct {
        u32  data_start, data_size;
        u8   checksum[16];
        u16  path_len;
        char path[path_len]; // NUL-term
    } entries[entry_count];
};
\end{lstlisting}

Walk the entries, \texttt{seek(data\_start)}, read \texttt{data\_size}, write
to \texttt{path}. Done: 1940 files fall out --- 1239 PNG + 246 JPG sprites, 363
XML layouts, 54 OGG tracks, and 29 Theora videos for the cutscenes. A Kaitai
Struct \texttt{.ksy} pins the whole format down in about 50 lines.

Extraction is free; \textit{repacking} is where the game
bites back. Every file carries a 16-byte checksum, and the TOC carries a global
one computed over \texttt{entry\_count} followed by each entry's
\texttt{(start, size, checksum, path\_len, path)}. On load, the engine
recomputes both and compares. Feed it a stock \texttt{md5()} and the result
\textit{looks} right --- correct 16 bytes --- yet the game rejects the archive.

The checksum \textbf{is} MD5, RFC~1321 to the letter: the same 64 sine
constants, the same rotation schedule, the same padding. Only the four
initialization words have been swapped:

\begin{lstlisting}[language=C]
//  stock MD5        Feeding Frenzy 2
A = 0x67452301  ->   0xfb132803
B = 0xefcdab89  ->   0x06e828f3
C = 0x98badcfe  ->   0xa12a051c
D = 0x10325476  ->   0x9bedf4fc
\end{lstlisting}

Four magic numbers, lifted straight out of \texttt{ff2.exe} in Binary Ninja.
Drop them into any MD5 implementation and the digests suddenly match the ones
stored in the original archive. That is the entire trick: a hash that is 99\%
standard and 1\% lie, betting you will reach for \texttt{hashlib} and never
notice the initial state is off.

With the real IV, a $\sim$100-line repacker rebuilds a byte-identical archive
from the extracted \texttt{resource/} tree --- recomputing every per-file digest
and the global one --- and the game loads it without complaint. Swap a sprite,
retranslate the XML, drop in your own fish: the tank is open.

\vspace{0.6em}
{\small Extractor, repacker, and format notes:\\
\url{https://github.com/xusheng6/feeding-frenzy-saf}}

\end{document}
